Software systems in safety-critical domains such as medical imaging and
industrial automation evolve continuously, making manually maintained
behavioural documentation increasingly difficult to keep up to date.
\sysml Interaction Models, structure and sequence diagrams that capture
the components involved in a scenario and the interactions between them,
are a natural target for automated reconstruction, but existing
reconstruction approaches rely on partial information sources and
therefore provide an incomplete view of system interactions: static
analysis over-approximates possible behaviour and cannot by itself
determine which interactions are exercised during a specific execution,
while execution traces provide execution-specific information but only
for behaviour exercised during the observed runs. This thesis
investigates: (RQ1) whether combining a static call graph with runtime
execution traces reconstructs scenario-specific Interaction Models more
accurately than either source alone; (RQ2) whether guiding trace
collection with the static call graph removes reconstruction noise
without losing scenario coverage; and (RQ3) how well a reconstructed
\sysml Interaction Model consolidates static and dynamic evidence into a
single artefact sufficient for system-behaviour comprehension.

\textbf{Methods:} A hybrid pipeline was built that synthesises Interaction
Models for event-driven \Cpp\ systems by combining a static call graph
with timestamped runtime traces collected by a separate pipeline. Static
and dynamic evidence were maintained as separate artefacts due to their
different evolution rates and storage characteristics; the call graph,
stored in a \neo property graph, contains only trace anchors, preserving
traceability without embedding large execution traces. The scenario-scoped
Cypher queries isolated the subset of static interactions associated with
a given scenario, where relevance was established through trace analysis
and corresponding static call-graph anchors. The pipeline was implemented using tools that, if required, can be invoked via twelve composable, natural-language-invocable agent skills. The tool was evaluated on \emph{CPSCore}, an open-source event-driven framework, across
four found-in-the-wild scenarios drawn from its existing test suite. Each
scenario was reconstructed under nine conditions (C1--C9), five of them
fully deterministic and reproducible (C2--C6): an \ac{LLM}-only baseline
(C1); static call-graph traversal (C2); unfiltered dynamic trace (C3);
guided instrumentation only (C4); guided instrumentation with a runtime
reachability filter (C5); full-trace instrumentation with the same filter
(C6); and three \ac{LLM}-augmented conditions (C7--C9) that post-process
C2, C3, and C4 candidates respectively. A
blinded six-rater expert panel compared diagrams reconstructed from
unfiltered trace output against diagrams reconstructed from trace output
scoped using static call-graph analysis of the scenario. The underlying
static approach was further tested in a practitioner session on a second,
larger industrial codebase to assess industrial usefulness beyond
reconstruction accuracy.

\textbf{Results:} Guided instrumentation with a Neo4j runtime reachability
filter (C5) achieved perfect F1 on both interaction-set and execution-order
metrics across all four scenarios (macro-F1~$=$~1.000\,/\,1.000). Full-trace
instrumentation with the same filter (C6) matched C5 exactly, showing the
filter is sufficient regardless of instrumentation strategy. Guided
instrumentation without the filter (C4) is the strongest Neo4j-free
deterministic baseline (macro-F1~$=$~0.914), while unfiltered dynamic tracing
(C3) was the weakest (0.530), dominated by intra-component false positives.
The \ac{LLM}-only baseline (C1) reached macro-F1~$=$~0.964, failing only on
S3 where a teardown-phase call appears in the source snippet. \ac{LLM}
augmentation of evidence priors (C7--C9) showed diminishing returns as the
prior improved: C7 substantially recovered C2's recall on three of four
scenarios; C8 partially denoised C3; C9 degraded C4's already-precise output
(0.783\,vs.\ C4's~0.914), confirming that \ac{LLM} post-processing is
counterproductive when the prior is already clean. Code-graph-guided
instrumentation eliminated all noise (50\%
fewer nodes, 70\% fewer edges) without losing scenario coverage, and was
rated higher by the expert panel in 47 of 54 comparisons. In a
scenario-comprehension task, no single evidence source alone, including a
reconstructed diagram, was sufficient: source code alone was weakest,
missing components and interactions that trace or diagram evidence made
explicit, while a diagram alone tied with runtime traces alone; only
combining source code with runtime traces reached perfect recall. In the
practitioner session, structural and dependency diagrams were rated
consistently useful, while sequence-diagram message labelling was
identified as the clearest area for improvement.

\textbf{Evaluation and Conclusion:} The three results converge on one
pattern: additional evidence improves reconstruction or comprehension
specifically when the primary evidence source is missing something, and
is otherwise redundant. The thesis contributes an empirical assessment of
hybrid static and dynamic analysis for reconstructing scenario-specific
\sysml diagrams. The results show when combining static and dynamic
evidence improves reconstruction, when scenario-guided instrumentation
reduces reconstruction noise without sacrificing coverage, and to what
extent reconstructed diagrams can replace the underlying evidence in
program-comprehension tasks. The work further provides an implementation
comprising a scenario-scoping method, a reversible instrumentation tool
(\texttt{csp\_matcher}), and an automated workflow based on
natural-language agent skills.
